\chapter{Representación de la Información}

Para que las funciones lógicas y los circuitos descritos en capítulos anteriores tengan utilidad práctica en el mundo de la ingeniería, necesitamos vincular los estados electromagnéticos abstractos (los ceros y unos del hardware) con entidades del mundo real: números naturales, enteros, reales, o caracteres de nuestro lenguaje escrito. Este capítulo formaliza el paso del sustrato físico-lógico hacia la semántica interpretativa.

\section{Sintaxis Lógica: Alfabetos y Lenguajes}

Antes de asignar significado, debemos definir rígidamente la estructura de lo que vamos a interpretar.
Un \textbf{alfabeto} $\Sigma$ es un conjunto finito, no vacío, de símbolos indivisibles. En electrónica digital elemental, nuestro alfabeto principal es el alfabeto booleano $\mathbb{B}_2 = \{0, 1\}.$ Sin embargo, la teoría es generalizable a cualquier alfabeto finito $\mathbb{D}_B = \{d_0, d_1, \dots, d_{B-1}\},$ donde $B$ es la cardinalidad o base del alfabeto.

Un \textbf{lenguaje formal} sobre un alfabeto $\Sigma$ es un conjunto de \textit{palabras}, donde una palabra se define como una secuencia (ristra) ordenada y finita de símbolos concatenados extraídos de $\Sigma.$ El conjunto de todas las palabras posibles de cualquier longitud, incluyendo la palabra vacía, se denota mediante la Clausura de Kleene $\Sigma^*.$

En el contexto de la arquitectura de computadores, el hardware está compuesto por buses y registros de un tamaño rígidamente predefinido. Por ello, no trabajamos con $\Sigma^*,$ sino que nos restringimos al subconjunto de palabras de \textbf{longitud constante} $n.$ 
Podemos visualizar formalmente una palabra $w$ de longitud $n$ como un vector perteneciente al espacio vectorial unidimensional o producto cartesiano $(\mathbb{D}_B)^n:$
\[ w = (w_{n-1}, w_{n-2}, \dots, w_1, w_0) \quad \text{donde } w_i \in \mathbb{D}_B \]
Al símbolo de índice $0$ ($w_0$) se le denomina dígito menos significativo (\textit{Least Significant Digit}, LSB en binario), y al símbolo de índice $n-1$ ($w_{n-1}$) dígito más significativo (\textit{Most Significant Digit}, MSB). Hasta este punto, una palabra es una pura estructura sintáctica sin valor intrínseco.

\section{Semántica: Números Naturales en Base $B$}

Para dotar de significado numérico a las palabras sintácticas, establecemos una aplicación semántica (función de evaluación) $V: (\mathbb{D}_B)^n \to \mathbb{N}.$
La convención universal para esta aplicación es la \textbf{notación posicional}, donde el aporte de cada símbolo al valor total está ponderado exponencialmente según su índice (posición) $i.$
Dado el alfabeto $\mathbb{D}_B$ con cardinal $B,$ asociamos a cada símbolo $d_k$ un valor primitivo $k.$
El valor natural total $V$ representado por la palabra $w = w_{n-1} \dots w_0$ es:
\[ V(w) = \sum_{i=0}^{n-1} V(w_i) \cdot B^i \]
Por ejemplo, en el alfabeto decimal ($B=10$), la palabra $402$ se evalúa como $4 \cdot 10^2 + 0 \cdot 10^1 + 2 \cdot 10^0 = 402.$ En el alfabeto binario ($B=2$), la palabra $1101_2$ evalúa a $1 \cdot 2^3 + 1 \cdot 2^2 + 0 \cdot 2^1 + 1 \cdot 2^0 = 13_{10}.$

Para un alfabeto arbitrario de base $B$ y una longitud de palabra de $n$ posiciones, el mínimo valor representable es $0$ (todas las posiciones en $d_0$) y el máximo valor representable está acotado algebraicamente por:
\[ V_{\max} = \sum_{i=0}^{n-1} (B-1) \cdot B^i = B^n - 1 \]
Toda operación aritmética cuyo resultado sobrepase $B^n - 1$ sufrirá pérdida de información o \textit{Desbordamiento (Overflow)}.

\section{Operaciones a Nivel de Palabra (Sintácticas)}

Previo a introducir semánticas más complejas como los números negativos, resulta imperativo definir matemáticamente dos operaciones unarias puramente sintácticas aplicables a cualquier palabra en un alfabeto de base $B.$

\subsection{Complementación a la Base menos 1 (Cb-1)}
Dada una palabra $w,$ el complemento a la base menos 1, que denotaremos provisoriamente como $C_{B-1}(w),$ se obtiene invirtiendo cada símbolo individualmente contra el valor máximo del alfabeto ($B-1$). 
\[ (C_{B-1}(w))_i = (B - 1) - w_i \]
En el caso particular del alfabeto binario ($B=2$), el valor $B-1 = 1,$ por lo que la operación se reduce a $(C_1(w))_i = 1 - w_i.$ Esta es estrictamente la definición de la compuerta lógica NOT ($\neg$). Por lo tanto, el Complemento a 1 de un número binario es simplemente su inversión bit a bit lógica.

\subsection{Complementación a la Base (Cb)}
El Complemento a la Base, denotado $C_B(w),$ se define matemáticamente partiendo del anterior y sumando $1$ a la palabra resultante:
\[ C_B(w) = C_{B-1}(w) + 1 \]
Algebraicamente, esto es equivalente a calcular $B^n - V(w).$ Dado que operamos en hardware de longitud fija $n,$ la suma genera frecuentemente un acarreo que escapa a la posición $n,$ el cual simplemente se descarta (operación módulo $B^n$).
En binario, esto corresponde al célebre Complemento a 2, que consiste en invertir todos los bits y sumar $1.$

\section{Semántica Avanzada: Enteros con Signo}

Para procesar números negativos (el conjunto $\mathbb{Z}$), la solución trivial de reservar un símbolo físico exclusivo tipo "$-$" es imposible, ya que el hardware solo dispone de su alfabeto base (ej. $0$ y $1$). Por ende, el signo debe ser codificado implícitamente empleando parte de la entropía de los símbolos existentes.

\subsection{Magnitud y Signo (M\&S)}
En este esquema, dividimos rígidamente la palabra $w$ de longitud $n$ en dos campos semánticos:
\begin{enumerate}
    \item El símbolo más significativo ($w_{n-1}$) se reserva exclusivamente para el signo. Por convención unánime en computación binaria, $0$ representa positivo y $1$ negativo.
    \item Los $n-1$ símbolos restantes se interpretan como la magnitud absoluta del número en notación posicional natural.
\end{enumerate}
El valor representado es:
\[ V(w) = (-1)^{w_{n-1}} \cdot \sum_{i=0}^{n-2} V(w_i) \cdot B^i \]
Este modelo presenta dos graves problemas de ingeniería: 
Primero, genera una doble representación del cero ($+0$ y $-0$), un desperdicio entrópico inaceptable.
Segundo, requiere de una Unidad Lógica Aritmética (ALU) notablemente compleja, ya que la suma y resta de variables requieren circuitos distintos dependientes del análisis de los signos previos.

\subsection{Complemento a la Base general (Cb)}
El estándar computacional moderno para representar enteros solventa los problemas anteriores asumiendo la semántica de la complementación a la base (Complemento a 2 en binario).
En este esquema, el dígito $w_{n-1}$ sigue decidiendo el signo, pero lo hace asignando un peso posicional estrictamente \textbf{negativo} al MSB. La ecuación de valor para un sistema binario en C2 es:
\[ V(w) = -w_{n-1} \cdot 2^{n-1} + \sum_{i=0}^{n-2} w_i \cdot 2^i \]
Generalizado a base $B,$ los números negativos se construyen calculando explícitamente el complemento a la base de la magnitud absoluta.
Las ventajas son abrumadoras:
\begin{itemize}
    \item Existe un único cero (el $-0$ colapsa tras sumar 1 y descartar el rebose).
    \item Rango asimétrico maximizado: $[-B^{n-1}, B^{n-1} - 1].$
    \item \textbf{Ceguera Aritmética:} La maravilla del Cb es que las restas $A - B$ se realizan como sumas $A + C_B(B).$ El hardware sumador (Full-Adder) no necesita saber si está sumando naturales, enteros positivos o negativos; el mismo circuito lógico produce el resultado binario correcto módulo $B^n.$
\end{itemize}

\subsection{Exceso a Bias (Exceso a $K$)}
Una alternativa para enteros es la semántica desplazada. En lugar de manipular el signo bit a bit, redefinimos el cero desplazándolo al centro del rango representable natural. Elegimos una constante de sesgo o \textit{Bias} $K.$
El valor representado $V$ es la evaluación natural estricta de la palabra binaria $V_{\text{nat}}(w)$ menos la constante $K:$
\[ V(w) = V_{\text{nat}}(w) - K \]
Convencionalmente, para $n$ dígitos, $K = B^{n-1}$ o bien $K = B^{n-1}-1.$ Este sistema garantiza que los números más negativos estén codificados sintácticamente con todos los ceros, y los más positivos con todos los unos. Aunque no es ideal para aritmética general (requiere ajustar el bias tras cada suma), es el sistema predilecto universal para comparar exponentes, dado que el hardware del comparador de magnitud binario estándar funciona de forma directa sin alteraciones.

\section{Representación Fraccionaria (Punto Fijo)}

Para representar los números Reales ($\mathbb{R}$), comenzamos asumiendo un "Punto Radix" o coma decimal inamovible (implícita por diseño hardware). 
Separamos la palabra de longitud total $N$ en dos tramos fijos: $n$ posiciones para la parte entera y $m$ posiciones para la parte fraccionaria ($N = n + m$).

Las ponderaciones posicionales decrecen, continuando a través del límite cero hacia potencias negativas de la base $B.$ 
La ecuación de un natural de punto fijo es la generalización natural:
\[ V(w) = \sum_{i=-m}^{n-1} V(w_i) \cdot B^i \]
Tanto la representación de Magnitud y Signo fraccionaria como la del Complemento a la Base fraccionario operan idénticamente al caso entero. En Complemento a 2 fraccionario, el bit de signo ocupa la posición ponderada a $-2^{n-1},$ y el resto de la estructura se mantiene, permitiendo a la ALU sumar decimales como si fueran enteros, simplemente ignorando lógicamente dónde ubicó el diseñador la coma de forma imaginaria.

\section{Representación en Punto Flotante}

El punto fijo es rígido: o perdemos exactitud en valores pequeños (poco tramo fraccionario) o perdemos alcance en los valores inmensos (poco tramo entero).
La notación científica solventa esto disociando el tamaño del número (el orden de magnitud) de sus cifras significativas. El estándar generalizado de punto flotante descompone la palabra binaria general en tres campos lógicos yuxtapuestos:
\begin{enumerate}
    \item \textbf{Signo (S):} Un solo bit (o equivalente en otra base) donde $0$ es positivo y $1$ negativo.
    \item \textbf{Exponente (E):} Una secuencia de $k$ posiciones tratadas siempre como un entero representado en \textbf{Exceso a Bias}. Este Bias permite manejar exponentes negativos.
    \item \textbf{Mantisa o Fracción (M):} Las $m$ cifras significativas del valor, operando estrictamente como un campo de coma fija fraccionario. 
\end{enumerate}

El formato impone una \textbf{Normalización}. En notación científica estándar, un número siempre se desliza para que la coma recaiga tras el primer dígito no nulo (ej. $5.32 \times 10^4$). En la base binaria ($B=2$), el único dígito distinto de $0$ es el $1.$ Esto nos regala un ahorro enorme: no necesitamos almacenar explícitamente el '1' inicial, ahorrando un bit completo. A este '1' no almacenado se le denomina \textbf{bit fantasma} o bit implícito.

La ecuación unificada de evaluación de un número de punto flotante normalizado en base binaria genérica es:
\[ V = (-1)^S \cdot (1 + M_{fracc}) \cdot 2^{E_{\text{nat}} - \text{Bias}} \]

\subsection{El Estándar IEEE-754}
El estándar real sobre el que funciona toda la electrónica y computación modernas para la precisión fraccionaria es el \textbf{IEEE-754}. Su encarnación más famosa son los tipos de datos nativos de punto flotante de los lenguajes de programación:

\begin{itemize}
    \item \textbf{Precisión Simple (float, 32 bits):} Un bit de signo, un exponente $E$ de 8 bits (Bias = $127$), y una mantisa $M$ de 23 bits fraccionarios.
    \item \textbf{Doble Precisión (double, 64 bits):} Un bit de signo, exponente de 11 bits (Bias = $1023$), y una descomunal mantisa de 52 bits para altísima precisión computacional.
\end{itemize}
Además, el IEEE-754 reserva combinaciones particulares del exponente y la mantisa para semánticas de excepción, como representar el Infinito ($\pm \infty$), el Cero ($\pm 0$ matemáticamente estricto) o entidades No Numéricas (\textit{Not a Number, NaN}), vitales cuando se intenta dividir por cero o hacer una raíz cuadrada de un número negativo.

\section{Códigos Alfanuméricos y Códigos Especiales}

La semántica posicional sirve para modelar la matemática. Pero para procesar cadenas de texto o interaccionar con entornos físicos, debemos establecer codificaciones puramente convencionales, diccionarios de traducción (Look-up tables).

\subsection{Códigos de Caracteres}
\begin{itemize}
    \item \textbf{ASCII:} El \textit{American Standard Code for Information Interchange} original codificaba 128 caracteres utilizando un alfabeto binario de longitud $n=7$ bits. Mapeaba los valores numéricos decimales del 0 al 127 contra los caracteres anglosajones y de control de los teletipos (por ejemplo, el 65 corresponde a la 'A' mayúscula y el 32 al espacio en blanco). Se extendió posteriormente a 8 bits (Añadiendo eñes, vocales acentuadas, etc).
    \item \textbf{Unicode:} Un inmenso consorcio internacional dedicado a codificar todos los sistemas de escritura humanos (desde el alfabeto latino, cirílico o chino, hasta símbolos matemáticos y emojis). Un identificador Unicode (Code Point) es un valor teórico abstracto, no una forma de escribir bits en memoria.
    \item \textbf{UTF (Unicode Transformation Format):} Son las formas reales de inyectar Unicode en memoria. 
    \begin{itemize}
        \item \textbf{UTF-8:} Estándar de longitud \textit{variable}. Mapea cada carácter usando de 1 a 4 bytes. Es universalmente retrocompatible con el ASCII básico.
        \item \textbf{UTF-16 y UTF-32:} Ocupan tamaños mínimos de 2 bytes (16 bits) o 4 bytes fijos (32 bits), útiles en bases de datos pesadas que prefieren ancho fijo para facilitar la indexación a costa de desperdiciar memoria.
    \end{itemize}
\end{itemize}

\subsection{Códigos de Distancia Unitaria (Código Gray)}
Existen dominios electromecánicos (como codificadores rotatorios o discos ópticos en brazos robóticos) donde el código posicional numérico puro de base 2 es desastroso. Si pasamos del número binario 3 ($011_2$) al número 4 ($100_2$), tres bits cambian de estado de forma completamente simultánea. Físicamente, un contacto eléctrico inevitablemente conmutará picosegundos antes que el otro, provocando que la máquina, durante una fracción minúscula de tiempo, lea un número intermedio basura disparando rutinas de fallo.

Para solucionar esto se inventaron los códigos continuos de distancia unitaria, el más famoso siendo el \textbf{Código Gray}.
En el código Gray, dos valores consecutivos cualquiera **siempre** difieren en exactamente un único bit (distancia de Hamming = $1$). Esto garantiza transiciones mecánicas inmaculadas sin lecturas esporádicas. Adicionalmente, esta propiedad de adyacencia de Gray es precisamente la infraestructura que organiza los ejes lógicos en los Mapas de Karnaugh que estudiamos en minimización.
